Vision-and-language navigation agents follow a written route instruction through a building using
only their own camera images, and an episode counts as a success only if the agent itself stops
near the goal. Running every saved checkpoint in a simulator is expensive, so a general-purpose
fine-tuning pipeline can instead keep the checkpoint with the lowest validation loss. For a
navigation policy this loss is teacher-forced: the model sees the expert's past frames and is scored
on whether it predicts the expert's next action. We fine-tune the Qwen3-VL-8B vision--language model
on Room-to-Room in Continuous Environments (R2R-CE), with and without class weighting, and run the
checkpoints this loss selects in the Habitat simulator, in buildings seen during training. None of
the selected checkpoints reaches the goal and stops there. Those trained without class weighting
never stop and behave like a baseline that always moves forward, and class weighting changes which
action the policy repeats. A preregistered follow-up executes every saved checkpoint, and no
checkpoint was shown to navigate better than the selected one. The validation loss gives no warning,
because stopping is rare among the decisions it averages, as earlier work has noted. Two offline
checks flag every run before any simulator test: recall on the \textsc{stop} action, and a pass/fail
test fixed before the first run. In this setting, a low validation loss is not evidence that a
policy can navigate. Code and run records are at
\href{https://github.com/vivekmattam02/Stop_Blind}{\texttt{github.com/vivekmattam02/Stop\_Blind}}.
